\section{بررسی اثر برش‌ها}
\label{sec:cutflow}
جدول~\ref{tab:cutflow} تعداد جت‌های باقی‌مانده را بعد از هر برش نشان می‌دهد. ترتیب برش‌ها این است: ابتدا \lr{$\pt>20$~GeV} (آستانه ذخیره‌سازی)، سپس \lr{$|\eta|<2$}، سپس \lr{$\pt>400$~GeV} و در پایان انتخاب منشأ (فقط جت کامل).
\begin{table}[!ht]
\centering
\caption{تعداد جت‌ها بعد از هر برش (سیگنال: \lr{zj}، \ttZ؛ پس‌زمینه: \bbbb، \ttbb).}
\label{tab:cutflow}
\begin{tabular}{lrrrr}
\toprule
برش & \lr{zj} & \ttZ & \bbbb & \ttbb \\
\midrule
\lr{$\pt>20$~GeV} & 74083 & 110926 & 90296 & 120574 \\
\lr{+ $|\eta|<2$} & 59905 & 90106 & 70128 & 96717 \\
\lr{+ $\pt>400$~GeV} & 28450 & 20871 & 24653 & 18500 \\
\lr{+ full} & 11733 & 14056 & 23071 & 12532 \\
\lr{full + partial} & 13197 & 14934 & 24375 & 13648 \\
\bottomrule
\end{tabular}
\end{table}

\widefig{origin_counts.pdf}{تعداد جت‌ها به تفکیک منشأ در هر نمونه، بدون برش، فقط با برش $|\eta|<2$، فقط با برش $\pt>400$ و با هر دو برش (محور عمودی لگاریتمی).}{fig:origcounts}
\widefig{origin_fractions.pdf}{سهم هر منشأ از کل جت‌های هر نمونه در همان چهار حالت.}{fig:origfrac}
شکل‌های~\ref{fig:origcounts} و~\ref{fig:origfrac} نشان می‌دهند هر برش چه سهمی از جت‌ها را حذف می‌کند و جت‌های باقی‌مانده از چه منشأیی‌اند. برش $\pt>400$ مؤثرترین برش است: تعداد جت‌ها را چند برابر کم می‌کند و بیشتر جت‌های نرم با منشأ پرتوهای اولیه، رویداد زمینه و تابش نرم را حذف می‌کند، در حالی که برش $|\eta|<2$ اثر کمتری دارد. در نتیجه سهم جت‌های دارای زوج \bb از تاپ‌نیامده در \bbbb از حدود ۴۲ درصد به حدود ۹۴ درصد و در \ttbb از حدود ۱۶ درصد به حدود ۶۸ درصد می‌رسد، و سهم جت‌های \Zbos{}$\to$\bb در \ttZ از حدود ۱۷ درصد به حدود ۶۷ درصد و در \lr{zj} از حدود ۲۳ درصد به حدود ۴۱ درصد. بخش بزرگ باقی‌مانده در \lr{zj} کوارک سبک سخت (حدود ۴۳ درصد) و در \ttZ و \ttbb جت‌هایی‌اند که \bq{} آنها از تاپ می‌آید (حدود ۲۳ درصد).

\label{sec:set1}
در این بخش توزیع هفت متغیر جت را برای هر چهار نمونه با هم نشان می‌دهیم: \pt، $\eta$، جرم، \lr{girth}، تعداد مؤلفه‌ها (\lr{nconst})، فاصله $\Delta R$ تا نزدیک‌ترین جت دیگر و تعداد جت در رویداد. هر متغیر در چهار مرحله برش رسم شده است تا ببینیم هر برش شکل توزیع‌ها را چطور تغییر می‌دهد. همه توزیع‌ها به مساحت واحد نرمال شده‌اند، مگر نمودارهای تعداد جت در رویداد که کسر رویدادها را نشان می‌دهند.

در هر یک از شکل‌های این مجموعه، چهار نمودار چهار مرحله برش را نشان می‌دهند: \lr{$\pt>20$}، \lr{+$|\eta|<2$}، \lr{+$\pt>400$}، و آخر داده نهایی.

\quadfig{s1_pT20_pt.pdf}{s1_eta_pt.pdf}{s1_pt400_pt.pdf}{s1_sel_pt.pdf}{توزیع \pt جت بعد از هر برش.}{fig:s1pt}
در مرحله اول (\lr{$\pt>20$}) توزیع \pt دو بخش دارد: یک بخش با افت تقریباً نمایی در \pt کم و یک برجستگی نزدیک ۴۰۰ تا ۵۰۰~\lr{GeV}. این برجستگی همان جت بوست‌شده‌ای است که برش پارتونی \lr{$\pt>400$} تولید می‌کند. بعد از برش \lr{$\pt>400$} فقط همین جت‌ها می‌مانند و توزیع یکنواخت می‌شود و تقریباً نمایی می‌افتد.

دو نکته درباره توزیع \pt جت‌ها لازم به ذکر است. نخست اینکه در \pt بالای ۷۰۰~\lr{GeV} جت‌های \ttZ و \ttbb نسبت به \Zjet و \bbbb کمی بیشترند، یعنی دم توزیع در دو فرایند حاوی تاپ سخت‌تر است. دوم اینکه افزایش ظاهری در آخرین \lr{bin} اثر نمایشی است، چون جت‌های بالای ۱۵۰۰~\lr{GeV} در آخرین \lr{bin} جمع می‌شوند.

دم سخت‌تر در فرایندهای حاوی تاپ می‌تواند چند دلیل داشته باشد که با نمودارهای این گزارش جور درمی‌آید. نخست، تفاوت از خود \Zbos{} شروع می‌شود: در خروجی \lr{MadGraph} سهم \Zbos{}هایی با $\pt>700$~\lr{GeV} در \ttZ حدود \lr{6.7} درصد و در \Zjet حدود \lr{5.6} درصد است (شکل~\ref{fig:lhe-prod})، با اینکه برش مولد در هر دو یکسان است ($\pt(\mathrm{Z})>400$~\lr{GeV}). این می‌تواند از انرژی بزرگ‌تر فرایند سخت در فرایندهای حاوی تاپ بیاید، چون میانه $\sqrt{\hat s}$ در \ttZ برابر ۱۴۹۲~\lr{GeV} در برابر ۱۰۴۸ در \Zjet و در \ttbb برابر ۱۵۱۲ در برابر ۱۱۰۶ در \bbbb است (بخش~\ref{sec:zorigin})، و انرژی بیشتر تکانه بیشتری برای \Zbos{} و سیستم همراهش باقی می‌گذارد. دوم، آلودگی ناشی از واپاشی تاپ است: در \ttZ حدود \lr{8.3} درصد از جت‌های کامل \Zbos{} دست‌کم یک \bq{} از تاپ دارند و تکانه این ذره‌های اضافه به جت افزوده می‌شود، در حالی که در \Zjet چنین منبعی وجود ندارد. سوم، قطبش \Zbos{} است: کسر طولی در \ttZ برابر $f_L\approx0.57$ و در \Zjet برابر \lr{0.07} است، و \Zbos{} طولی تکانه را متقارن‌تر بین دو \bq{} پخش می‌کند و آنها را به هم نزدیک‌تر می‌کند، پس \Zbos{} بهتر در مخروط $R=0.8$ جا می‌گیرد و تکانه کمتری بیرون از جت می‌ماند. نشانه این اثر سهم جت‌های جزئی است که در \Zjet حدود ۱۱ درصد و در \ttZ حدود ۶ درصد است؛ این جت‌ها \Zbos{} را در مخروط دارند ولی کمتر از دو \bq{} و در نتیجه بخشی از تکانه \Zbos{} را از دست داده‌اند.

برای \ttbb در برابر \bbbb باید به برش‌های مولد هم توجه کرد، چون با هم فرق دارند: در \bbbb هر دو زوج باید $\pt>400$~\lr{GeV} و $\Delta R<0.8$ داشته باشند، ولی در \ttbb فقط یک زوج با $\pt>400$~\lr{GeV} و $\Delta R<0.8$. ساختار پس‌زدن هم فرق دارد: بعد از برش‌های $|\eta|<2$ و $\pt>400$ حدود ۵۱ درصد رویدادهای \Zjet و ۴۰ درصد \bbbb دو جت دارند و جت دیگر پشت‌به‌پشت جت منتخب است (قله $\Delta R\approx\pi$ در شکل~\ref{fig:s1dr})، ولی در \ttZ و \ttbb این سهم فقط حدود ۲۲ و ۱۸ درصد است (شکل~\ref{fig:s1njets}).

\quadfig{s1_pT20_eta.pdf}{s1_eta_eta.pdf}{s1_pt400_eta.pdf}{s1_sel_eta.pdf}{توزیع $\eta$ جت بعد از هر برش.}{fig:s1eta}
بدون برش \lr{$\eta$}، توزیع‌ها تا $|\eta|\approx5$ ادامه دارند و برای چهار نمونه تقریباً یکسان‌اند. بعد از برش \lr{$\pt>400$} جت‌ها مرکزی‌تر می‌شوند. در شکل نهایی \ttZ مرکزی‌ترین نمونه است (قله بلندتر در $\eta=0$)، \Zjet تخت‌ترین، و \bbbb و \ttbb بین آنها قرار دارند. این تفاوت در خود \Zbos{} هم هست و علتش آستانه انرژی فرایند سخت است. اگر دو پارتون ورودی کسرهای $x_1$ و $x_2$ از تکانه پروتون را داشته باشند، انرژی مرکز جرم فرایند سخت $\hat s=x_1x_2S$ و تندی سیستم $y=\frac12\ln(x_1/x_2)$ است، که $S=(13~\mathrm{TeV})^2$ مربع انرژی مرکز جرم برخورد دو پروتون است. از این دو رابطه $x_{1,2}=\sqrt{\hat s/S}\,e^{\pm y}$ به دست می‌آید و چون $x_{1,2}\le1$ است، همیشه $|y|\le\ln(\sqrt{S}/\sqrt{\hat s})$ است. پس هر چه $\sqrt{\hat s}$ بزرگ‌تر باشد بازه مجاز $y$ باریک‌تر می‌شود و رویدادها به مرکز نزدیک‌تر می‌شوند. \Zbos{} و جت مقابلش در جهت عرضی تعادل دارند و $\eta$ جت تندی همین سیستم را دنبال می‌کند: برای ذره‌ای که جرمش در برابر تکانه‌اش ناچیز است، یعنی سرعتش به سرعت نور نزدیک است، تندی $y$ تقریباً با شبه‌سرعت $\eta$ برابر است، و جت‌های ما با $\pt>400$~\lr{GeV} در همین وضعیت‌اند. پس سیستم باریک‌تر یعنی جت مرکزی‌تر. در \ttZ علاوه بر \Zbos{} باید دو تاپ هم ساخته شوند، پس کمینه $\sqrt{\hat s}$ بالاتر است: با $\pt(\mathrm{Z})>400$~\lr{GeV} این کمینه تقریباً ۸۱۰~\lr{GeV} برای \Zjet و ۹۴۰~\lr{GeV} برای \ttZ است. داده همین را نشان می‌دهد: میانه $\sqrt{\hat s}$ برابر ۱۴۹۲ در \ttZ و ۱۰۴۸ در \Zjet و ۱۵۱۲ در \ttbb و ۱۱۰۶ در \bbbb است (بر حسب \lr{GeV})، و میانگین $|y|$ سیستم $0.55$ در برابر $0.80$ و $0.52$ در برابر $0.61$ است.

\quadfig{s1_pT20_mass.pdf}{s1_eta_mass.pdf}{s1_pt400_mass.pdf}{s1_sel_mass.pdf}{توزیع جرم جت بعد از هر برش.}{fig:s1mass}
در مرحله \lr{$\pt>20$} بیشتر جت‌ها جرم کوچک دارند (جت‌های سبک و نرم). با برش \lr{$\pt>400$} و سپس انتخاب منشأ، در دو نمونه سیگنال قله‌ای تیز نزدیک جرم \Zbos{} (حدود ۸۰ تا ۱۰۰~\lr{GeV}) ظاهر می‌شود. در شکل نهایی بیشتر جت‌های \Zjet و \ttZ در این قله‌اند. نمونه‌های پس‌زمینه قله مشخصی ندارند: توزیعشان پهن است و نزدیک ۶۰ تا ۷۰~\lr{GeV} یک برجستگی کم‌ارتفاع دارد و دم بلند به جرم‌های بالا.

برجستگی پس‌زمینه‌ها در حدود ۶۰ تا ۷۰~\lr{GeV} با سرکوب سودکوف قابل فهم است. جرم جت از تابش گلوئون می‌آید: اگر ذره‌ای نرم با سهم انرژی $z$ و زاویه $\theta$ از محور جت تابش شود، جرمی با $m^2\approx z\,\theta^2\,\pt^2$ می‌دهد. در تقریب نرم و هم‌خط، احتمال این تابش $\frac{\alpha_s C}{\pi}\frac{d\theta^2}{\theta^2}\frac{dz}{z}$ است که $C$ ضریب رنگ است ($C_A=3$ برای گلوئون و $C_F=4/3$ برای کوارک). با انتگرال‌گیری روی $z$ و روی $\theta<R$، احتمال داشتن دست‌کم یک تابش با جرم بیشتر از $m$ برابر $\frac{\alpha_s C}{2\pi}L^2$ می‌شود که $L=\ln\frac{(R\pt)^2}{m^2}$ است، و احتمال اینکه هیچ تابشی بالاتر از $m$ نباشد، یعنی سهم جت‌هایی با جرم کمتر از $m$، برابر است با
\[
\Sigma(m)=\exp\!\Big(-\frac{\alpha_s C}{2\pi}\,L^2\Big).
\]
این همان سرکوب سودکوف است: هر چه $m$ کوچک‌تر باشد $L$ بزرگ‌تر است و $\Sigma$ شدیدتر به صفر می‌رود، پس جت بدون هیچ تابشی و با جرم نزدیک صفر بسیار نادر است. با $a=\alpha_s C/\pi$ مشتق آن توزیع جرم را می‌دهد: $d\sigma/d\ln m\propto a\,L\,e^{-aL^2/2}$ که در $L=1/\sqrt{a}$ بیشینه دارد، یعنی در $m_\mathrm{peak}=R\,\pt\,e^{-1/(2\sqrt{a})}$. چون نمودارهای ما چگالی بر حسب خود $m$ هستند، یک عامل $1/m$ هم وارد می‌شود و بیشینه از $aL^2-L/2-1=0$ به دست می‌آید که قله را به جرم‌های کوچک‌تر می‌برد. برای $\pt\approx450$~\lr{GeV} و $R=0.8$ (یعنی $R\pt=360$~\lr{GeV}) و $0.10<\alpha_s<0.14$، برای گلوئون قله $d\sigma/d\ln m$ بین ۷۱ و ۹۲~\lr{GeV} و قله $d\sigma/dm$ بین ۱۲ و ۲۷~\lr{GeV} است؛ برای کوارک این مقدارها ۳۲ تا ۴۶~\lr{GeV} و کمتر از ۳~\lr{GeV} هستند. بنابراین مرتبه بزرگی قله، چند ده \lr{GeV}، با این برآورد سازگار است و نشان می‌دهد جرم جت \lr{QCD} به‌طور طبیعی به صفر میل نمی‌کند. \textcolor{red}{ولی جای دقیق قله (حدود ۶۰ تا ۷۰~\lr{GeV} در $d\sigma/dm$) را این برآورد تک‌لگاریتمی با ضریب $\alpha_s$ ثابت بازتولید نمی‌کند و به $\alpha_s$ و به متغیر رسم‌شده حساس است؛ نقش اثرهای دیگر مثل شکافت سخت، جرم \bq{}، هادرونیزاسیون و تفکیک‌پذیری آشکارساز را بررسی نکرده‌ایم.}

\quadfig{s1_pT20_girth.pdf}{s1_eta_girth.pdf}{s1_pt400_girth.pdf}{s1_sel_girth.pdf}{توزیع \lr{girth} بعد از هر برش.}{fig:s1girth}
بعد از برش‌ها، سیگنال در \lr{girth} حدود \lr{0.15} تا \lr{0.2} قله دارد (جت دو‌هسته‌ای که انرژی از محور فاصله دارد)، ولی پس‌زمینه در \lr{0.05} تا \lr{0.1} قله دارد و دم بلندتر دارد. دو نمونه سیگنال و دو نمونه پس‌زمینه با هم هماهنگ‌اند.

\quadfig{s1_pT20_nconst.pdf}{s1_eta_nconst.pdf}{s1_pt400_nconst.pdf}{s1_sel_nconst.pdf}{توزیع تعداد ذرات تشکیل‌دهنده جت (\lr{nconst}) بعد از هر برش.}{fig:s1nconst}
در مرحله اول، قله نزدیک ۱۰ ذره تشکیل‌دهنده از جت‌های نرم می‌آید. بعد از برش \lr{$\pt>400$} توزیع به قله‌ای در حدود ۴۰ تا ۵۰ ذره تشکیل‌دهنده تغییر می‌کند. در داده نهایی سیگنال کمی ذره تشکیل‌دهنده کمتر از پس‌زمینه دارد.

\quadfig{s1_pT20_dr.pdf}{s1_eta_dr.pdf}{s1_pt400_dr.pdf}{s1_sel_dr.pdf}{توزیع فاصله $\Delta R$ جت تا جت دیگر. در سه مرحله اول فاصله جت اول و دوم و در مرحله نهایی فاصله جت منتخب تا نزدیک‌ترین جت دیگر با \lr{$\pt>20$} نشان داده شده است.}{fig:s1dr}
در سه مرحله اول، $\Delta R$ در هر چهار نمونه نزدیک $\pi\approx3.1$ قله دارد، یعنی دو جت سخت معمولاً پشت‌به‌پشت‌اند. سمت راست قله ($\Delta R>\pi$) فقط از اختلاف $\eta$ می‌آید، چون اختلاف زاویه سمتی از $\pi$ بیشتر نمی‌شود. در داده نهایی متغیر فرق می‌کند و فاصله جت منتخب تا نزدیک‌ترین جت دیگر را نشان می‌دهد؛ قله پشت‌به‌پشت اینجا فقط در \Zjet و \bbbb دیده می‌شود و در \ttZ و \ttbb نیست. این با تعداد جت‌های رویداد سازگار است: با $\pt>20$ به‌طور میانگین \lr{3.7} جت در هر رویداد \Zjet، \lr{4.5} در \bbbb، \lr{5.5} در \ttZ و \lr{6.0} در \ttbb هست، و بعد از برش‌های $|\eta|<2$ و $\pt>400$ حدود ۵۱ درصد رویدادهای \Zjet و ۴۰ درصد \bbbb دقیقاً دو جت دارند، در حالی که در \ttZ و \ttbb این سهم حدود ۲۲ و ۱۸ درصد است (شکل~\ref{fig:s1njets}). یک توضیح ممکن این است که در رویدادهای حاوی تاپ، تکانه مقابل بین چند جت کم‌تکانه‌تر پخش می‌شود.

\quadfig{s1_pT20_njets.pdf}{s1_eta_njets.pdf}{s1_pt400_njets.pdf}{s1_sel_njets.pdf}{تعداد جت در هر رویداد بعد از هر برش.}{fig:s1njets}
تعداد جت‌ها در هر رویداد (شکل~\ref{fig:s1njets}) بدون هیچ آستانه‌ای حدود ۳۳ تا ۳۵ است (بازه ۶ تا ۵۳) و چهار نمونه در آن کاملاً شبیه‌اند. با $\pt>20$ این عدد به ۲ تا حدود ۱۴ می‌رسد و میانگینش \lr{3.7} در \Zjet، \lr{4.5} در \bbbb، \lr{5.5} در \ttZ و \lr{6.0} در \ttbb است؛ کمترین در \Zjet و بیشترین در \ttbb است، که می‌تواند از جت‌های اضافه واپاشی تاپ بیاید. بعد از برش‌های $|\eta|<2$ و $\pt>400$ حدود ۵۱ درصد رویدادهای \Zjet و ۴۰ درصد \bbbb دقیقاً دو جت دارند، در حالی که در \ttZ و \ttbb این سهم حدود ۲۲ و ۱۸ درصد است. در داده نهایی علاوه بر این برش‌ها شرط منشأ جت هم هست، پس در هر رویداد حداکثر یک جت می‌ماند، به‌جز \bbbb که دو زوج \bq{} دارد: در \Zjet حدود ۳۴ درصد رویدادها صفر جت و ۶۶ درصد یک جت دارند، در \ttZ حدود ۲۵ و ۷۵ درصد، در \ttbb حدود ۳۳ درصد صفر، ۶۶ درصد یک و ۱ درصد دو جت، و در \bbbb حدود ۱۸ درصد صفر، ۴۲ درصد یک و ۴۰ درصد دو جت. میانگین کمتر از یک جت در رویداد در \Zjet، \ttZ و \ttbb طبیعی است، چون بعضی رویدادها بعد از برش‌ها هیچ جت بوست‌شده‌ای ندارند. پس تفاوت نمونه‌ها از برش‌های سخت می‌آید، نه از رفتار کلی رویداد.
